\documentclass[12pt]{article}
\usepackage[utf8]{inputenc}
\usepackage[T1]{fontenc}
\usepackage{hyperref}
\title{Quioquitani-Quier\'i Zapotec Language Resources for K-12 Curriculum Development}
\author{Compiled by Claude Sonnet 5.5 (claude-sonnet-5-5) for Andr\'as Kornai}
\date{}
\begin{document}
\maketitle

\section{Introduction}
Quioquitani-Quier\'i Zapotec (ISO 639-3 ztq, Glottocode quio1241) is an Otomanguean language of Oaxaca, Mexico, spoken around Santa Catarina Quioquitani and Santa Catalina Quier\'i. Estimates of speakers range from about 4,000 to about 6,300. The community name for the language is given as Tiits S\"e in some sources. The searches found little beyond scripture-related and reference-catalogue material. No school primers, textbooks or dictionaries turned up.

\section{Primary Educational Resources}
No primers, textbooks or school curricula in this language were found in the searches. The only educational-type material found is the Bible and Bible-teaching audio listed below.

\section{Linguistic and Reference Materials}
\begin{itemize}
\item Wikitongues language page (Latin script, Otomanguean, developing status): \url{https://wikitongues.org/languages/ztq/}
\item ASJP wordlist entry: \url{https://asjp.clld.org/languages/QUIOQUITANI_QUIERI_ZAPOTEC}
\item Numeral-system description (Max Planck channumerals database): \url{https://lingweb.eva.mpg.de/channumerals/Zapotec-Quioquitani-Quieri.htm}
\end{itemize}

\section{Religious and Community Resources}
\begin{itemize}
\item Scripture Earth page with the New Testament (text, audio, app): \url{https://www.scriptureearth.org/00eng.php?iso=ztq}
\item eBible entry for the New Testament: \url{https://ebible.org/Bible/details.php?id=ztqNT}
\item Joshua Project people-group profile: \url{https://joshuaproject.net/people_groups/print/14736/MX}
\end{itemize}

\section{Digital and Audio Resources}
\begin{itemize}
\item Global Recordings Network audio and slideshow recordings: \url{https://globalrecordings.net/en/language/ztq}
\end{itemize}

\section{Curriculum Development Context}
The language appears to be used at home, with Spanish used for formal education. A New Testament in Latin script exists, which provides a basis for orthography and reading texts. Any K-12 curriculum would need new materials developed with the community. No existing school programme was found.
\end{document}